% Part: conditional-logics
% Chapter: minimal-change-semantics
% Section: sphere-models

\documentclass[../../../include/open-logic-section]{subfiles}

\begin{document}

\olfileid{con}{min}{sph}

\olsection{Model Bola}

Salah siji cara menehi semantik formal kanggo kontrafaktual
yaiku ngowahi panjelasan informal Lewis dadi struktur matematika.
Bola-bola ing saubenge donya~$w$ banjur minangka himpunan donya.
Amarga bola-bola iku mapan ing njero siji lan sijine, himpunan
donya ing saubenge~$w$ kudu diurutake kanthi linear dening
relasi himpunan bagean.

\begin{defn}
  \emph{Model bola} iku tripel~$\mModel{M} = \tuple{W, O, V}$,
  kanthi $W$ minangka himpunan donya kang ora kosong,
  $V\colon \PVar \to \Pow{W}$ minangka valuasi, lan
  $O\colon W \to \Pow{\Pow{W}}$ masangake saben donya~$w$
  karo \emph{sistem bola}~$O_w$. Kanggo saben $w$, $O_w$ iku
  himpunan kang anggotane himpunan donya lan kudu nyukupi:
  \begin{enumerate}
  \item $O_w$ \emph{dipusatake} ing~$w$: $\{w\} \in O_w$.
  \item $O_w$ \emph{susun ing njero siji lan sijine}: saben
    $S_1$, $S_2 \in O_w$, $S_1 \subseteq S_2$ utawa
    $S_2 \subseteq S_1$, yaiku $O_w$ diurutake kanthi linear
    dening~$\subseteq$.
  \item $O_w$ tutup tumrap gabungan saka kulawarga bola kang ora kosong.
  \item $O_w$ tutup tumrap irisan saka kulawarga bola kang ora kosong.
  \end{enumerate}
\end{defn}

Intuisi kang ndhasari $O_w$ yaiku donya-donya ``ing saubenge'' $w$
diperang dadi lapisan miturut adoh-cedhake saka~$w$. Bola ora
kosong kang paling njero mung ngemot $w$ dhewe, yaiku
himpunan~$\{w\}$: $w$ luwih cedhak karo~$w$ tinimbang donya-donya
liyane ing bola liyane. Yen $S \subsetneq S'$, mula donya-donya ing
$S' \setminus S$ luwih adoh saka $w$ tinimbang donya-donya
ing~$S$: $S' \setminus S$ iku ``lapisan'' antarane $S$ lan
donya-donya ing njaba~$S'$. Mligine, bola kudu dimangerteni
minangka ngemot kabeh donya ing njero lumahing njabane,
dudu mung lapisan dhewe-dhewe.

\begin{figure}
\begin{center}
\begin{tikzpicture}[layerwidth=1.5,scale=.6]\tiny
  \spheresystem{3}
  \propositionintersect{0}{3}{60}{6}
  \path (0,0) node[world] {$w$};
  \spherepos{-30}{2}{node[world] {$w_2$}}
  \spherepos{220}{2}{node[world] {$w_3$}}
  \spherepos{90}{2}{node[world] {$w_1$}}
  \spherepos[shift={(.1,0)}]{10}{3}{node[world] {$w_5$}}
  \spherepos[shift={(.1,0)}]{-10}{3}{node[world] {$w_6$}}
  \spherepos{225}{3}{node[world] {$w_4$}}
  \spherepos{20}{4}{node[world] {$w_7$}}
  \path (5,0) node[left] {$p$};
\end{tikzpicture}
\caption{Diagram model bola}
\ollabel{fig:sphere-model}
\end{center}
\end{figure}

Diagram ing \olref{fig:sphere-model} cocog karo model bola kanthi
$W = \{w, w_1, \dots, w_7\}$, $V(p) = \{w_5, w_6, w_7\}$.
Bola paling njero yaiku $S_1 = \{w\}$. Donya-donya kang paling
cedhak karo $w$ yaiku $w_1, w_2, w_3$, mula bola kang sabanjure
luwih gedhe yaiku $S_2 = \{w, w_1, w_2, w_3\}$. Donya-donya
kang luwih adoh yaiku $w_4$, $w_5$, $w_6$, mula bola paling
njaba yaiku $S_3 = \{w, w_1, \dots, w_6\}$. Sistem bola ing
saubenge $w$ yaiku $O_w = \{S_1, S_2, S_3\}$. Donya~$w_7$
ora ana ing bola apa wae ing saubenge~$w$. Donya-donya kang
paling cedhak lan ing kono $p$~bener yaiku $w_5$ lan~$w_6$,
mula bola paling cilik kang ngakoni $p$ yaiku~$S_3$.

Kanggo nemtokake pemenuhan formula~$!A$ ing donya~$w$ ing
model bola~$\mModel M$, $\mSat{M}{!A}[w]$, kita ngluwihi
definisi kanggo !!{formula}s modal kanthi nambah klausa
kanggo $!B \cif !C$:

\begin{defn}
  $\mSat{M}{!B \cif !C}[w]$ yen lan mung yen salah siji iki lumaku:
  \begin{enumerate}
  \item\ollabel{sphere-vac} Kanggo kabeh $u \in \bigcup O_w$,
    $\mSat/{M}{!B}[u]$, utawa
  \item\ollabel{sphere-nonvac} Kanggo sawenehing $S \in O_w$,
    \begin{enumerate}
      \item $\mSat{M}{!B}[u]$ kanggo sawenehing $u \in S$, lan
      \item kanggo kabeh $v \in S$, salah siji
        $\mSat/{M}{!B}[v]$ utawa $\mSat{M}{!C}[v]$ lumaku.
    \end{enumerate}
  \end{enumerate}
\end{defn}

Miturut definisi iki, $\mSat{M}{!B \cif !C}[w]$ yen lan mung
yen anteseden~$!B$ luput ing kabeh donya ing bola-bola saubenge~$w$,
utawa ana bola~$S$ kang ing kono $!B$ bener lan kondisional
material $!B \lif !C$ bener ing kabeh donya ing bola kang
``ngakoni $!B$'' iku. Gatekna yen definisi ora mbutuhake
$S$ minangka bola kang ngakoni $!B$ lan \emph{paling njero},
sanajan saka panjelasan intuitif kita bisa uga ngarepake
mangkono. Nanging yen syarat ing~\olref{sphere-nonvac}
dipenuhi dening sawenehing bola~$S$, syarat iku uga dipenuhi
dening kabeh bola ing njero~$S$ kang isih ngakoni anteseden.
Mula, mligine, syarat iku dipenuhi ing bola kang ngakoni
anteseden lan paling njero, yen bola kaya mangkono ana.

Gatekna uga yen definisi model bola ora mbutuhake yen
\emph{ana} bola kang ngakoni $!B$ lan paling njero: kita bisa
nduweni runtutan tanpa wates $S_1 \supsetneq S_2 \supsetneq
\dots \supsetneq \{w\}$ saka bola-bola kang ngakoni $!B$,
kanthi irisan kabeh bola ing runtutan iku mung donya pusat,
lan anteseden luput ing donya pusat. Ing kasus iki ora ana
bola kang ngakoni $!B$ lan paling njero. Banjur
$\mSat{M}{!B \cif !C}[w]$ yen lan mung yen $!B \lif !C$
lumaku ing kabeh donya ing bola-bola $S_i$, $S_{i+1}$,
\dots, kanggo sawenehing~$i$.

\end{document}
